\section{习题}

以下三十题按八讲分组，题号连续。每题先给题目，再给一句提示；提示只指出第一步，完整论证请自己补全。标 $\dagger$ 的题需要回到原文相应小节查阅细节。

\subsubsection*{第 1 讲\quad 舞台}

\begin{enumerate}[label=\arabic*.,start=1,leftmargin=2.4em,itemsep=6pt]
\item 设 $C$ 是亏格 $g\ge2$ 的紧 Riemann 面，$G$ 是单连通的单群。证明 Hitchin 模空间 $\MH$ 的复维数为 $(2g-2)\dim G$。
\par\smallskip\noindent\emph{提示：}$\dim\mathrm{Bun}_G=(g-1)\dim G$；$\varphi\in H^0(C,\mathrm{ad}E\otimes K_C)$，对稳定丛用 Riemann--Roch 与 Serre 对偶把 $h^0(\mathrm{ad}E\otimes K_C)$ 算成 $(g-1)\dim G$。

\item 取 $G=U(1)$。证明 $\MH\cong T^*\mathrm{Jac}(C)$，写出 Hitchin 纤维化及其纤维，并说明为什么上一题的维数公式对含 $U(1)$ 因子的群不适用。
\par\smallskip\noindent\emph{提示：}$\mathrm{Bun}_{U(1)}=\mathrm{Jac}(C)$，而 $H^0(C,K_C)\cong\CC^g$；公式失效的原因是 $H^0(C,\mathcal O)\neq0$ 使 $\chi(\mathrm{ad}E)$ 多出中心项。

\item 说明 Casimir 多项式的系数给出 Hitchin 纤维化 $\MH\to B=\bigoplus_iH^0(C,K_C^{d_i})$，并证明 $\dim B=(g-1)\dim G$。
\par\smallskip\noindent\emph{提示：}$\varphi\mapsto(\Tr\varphi^2,\Tr\varphi^3,\dots)$ 在纤维上取常值；再用 $\dim H^0(C,K_C^{d})=(2d-1)(g-1)$ 与 $\sum_i(2d_i-1)=\dim G$。

\item 把 $\MH$ 写成无穷维超凯勒商：写出平坦联络与 Higgs 场空间上的三个矩映射，并验证三个水平条件合起来就是 Hitchin 方程。
\par\smallskip\noindent\emph{提示：}三个条件分别是 $F_A+[\varphi,\varphi^*]=0$、$\bar\partial_A\varphi=0$ 以及后者的共轭；商去的群是规范群。
\end{enumerate}

\subsubsection*{第 2 讲\quad 4d 主角}

\begin{enumerate}[label=\arabic*.,start=5,leftmargin=2.4em,itemsep=6pt]
\item 证明 $\Ldual{SU(N)}=PSU(N)$，且 $\Ldual{PSU(N)}=SU(N)$。
\par\smallskip\noindent\emph{提示：}对偶群把根系与余根系互换；$SU(N)$ 的中心 $\ZZ_N$ 在对偶群里变成基本群 $\ZZ_N$。

\item 写出 $SU(2)$ 与 $SO(3)$ 的权格与余权格，并验证 GNO 对偶把二者互换。
\par\smallskip\noindent\emph{提示：}$SU(2)$ 的权格是 $(1/2)\ZZ$，$SO(3)$ 的权格是 $\ZZ$；磁荷所取的余权格恰是权格的对偶。

\item 由 $\tau=\theta/2\pi+4\pi\ii/e^2$ 读出 $\Re\tau=\theta/2\pi$、$\Im\tau=4\pi/e^2$，说明 $T:\tau\mapsto\tau+1$ 等价于 $\theta\mapsto\theta+2\pi$，并解释为什么 $\theta$ 的周期必须是 $2\pi$。
\par\smallskip\noindent\emph{提示：}瞬子数取整数值，$\theta$ 的平移只改变 $e^{\ii\theta n}$ 的相位；$T$ 变换就是这一平移。

\item 用 $S=\begin{pmatrix}0&1/\sqrt{n_g}\\-\sqrt{n_g}&0\end{pmatrix}$ 验证 $S:\tau\mapsto-1/(n_g\tau)$，并说明为什么 $n_g=2,3$ 时 $S$ 与 $T$ 生成的是 $SL(2,\RR)$ 的离散子群而不是整个 $SL(2,\ZZ)$。
\par\smallskip\noindent\emph{提示：}直接做矩阵乘法；$n_g$ 不是平方数时 $1/\sqrt{n_g}$ 无理，故 $S\notin SL(2,\ZZ)$。
\end{enumerate}

\subsubsection*{第 3 讲\quad 扭化}

\begin{enumerate}[label=\arabic*.,start=9,leftmargin=2.4em,itemsep=6pt]
\item 由 $\Psi=(\tau t^2+\bar\tau)/(1+t^2)$ 验证：$t=\pm1$ 时 $\Psi=\Re\tau$ 为实；$t=\pm\ii$ 时 $\Psi=\infty$ 且与 $\tau$ 无关；$|t|=1$ 时 $\Psi$ 为实。
\par\smallskip\noindent\emph{提示：}$t=\pm\ii$ 时分母为零；$|t|=1$ 时 $1/t=\bar t$，故 $(t-1/t)/(t+1/t)$ 为纯虚。

\item 设 $S$ 对 $(\tau,t)$ 的作用为 $\tau'=-1/(n_g\tau)$、$t'^2=t^2\tau/\bar\tau$。由 $\Psi$ 的定义验证 $\Psi'=-1/(n_g\Psi)$。
\par\smallskip\noindent\emph{提示：}把 $\Psi'$ 通分后代入 $\tau\bar\tau=|\tau|^2$，并注意 $t'^2\tau'$ 与 $t'^2\bar\tau'$ 的结构。

\item 固定 $\Psi$ 与 $\Re\tau$，令 $\Im\tau\to\infty$，证明 $t^2\to1$，即弱耦合极限对应 $t\to\pm1$。
\par\smallskip\noindent\emph{提示：}$t^2=(\bar\tau-\Psi)/(\Psi-\tau)$；写 $\tau=\Re\tau+\ii\Im\tau$，比较分子分母的最高阶项。

\item 说明为什么 $16$ 个超荷在 $Spin(4)_{space}\times Spin(6)_R$ 下分解、取对角 $Spin(4)$ 后恰有两个不变超荷，并解释为什么扭化族的参数是 $\PP^1$。
\par\smallskip\noindent\emph{提示：}$16=(2,1,4)+(1,2,\bar4)$；把 $Spin(6)$ 的 $4$ 限制到 $Spin(4)$ 后两条链各留下一个单态，两个超荷的复线性组合只差一个整体标度。
\end{enumerate}

\subsubsection*{第 4 讲\quad 紧化与镜像}

\begin{enumerate}[label=\arabic*.,start=13,leftmargin=2.4em,itemsep=6pt]
\item 解释为什么四维理论沿 $C$ 紧化、只保留零模后，二维理论是靶空间为 $\MH$ 的 sigma 模型。
\par\smallskip\noindent\emph{提示：}零模由平坦联络加全纯 Higgs 场刻画，而 BPS 方程正是 Hitchin 方程；sigma 模型的靶空间就是这些模构成的模空间。

\item 在 $\tau$ 为纯虚（$\theta=0$）时，验证 $t=\pm\ii$ 对应 $(w_+,w_-)=(-\ii,-\ii)$、因而是 $\MH(G,C)$ 上复结构 $J$ 的 B-模型；再施加 $S$ 对偶得到 $t=\pm1$ 与 $\Psi=0$，对应 $\MH(\Ldual G,C)$ 上复结构 $K$ 的 A-模型。
\par\smallskip\noindent\emph{提示：}判据是 $J_+=J_-$ 给 B-模型、$J_+=-J_-$ 给 A-模型；$S$ 把 $t=\pm\ii$ 送到 $t=\pm1$。

\item 验证 $S$-对偶把支撑在单点的 zerobrane（型 $(B,B,B)$）送到型 $(B,A,A)$ 的膜，并说明为什么后者的支撑必须是某条 Hitchin 纤维。
\par\smallskip\noindent\emph{提示：}查 Table 2 的变换 $(B,B,B)\mapsto(B,A,A)$；A-膜至少是中维子流形，而 zerobrane 的像维数不超过中维，故恰好落在 Hitchin 纤维上。

\item 说明 Hitchin 纤维 $F$ 与对偶群一侧的 Hitchin 纤维互为对偶环面，并解释 $S$-对偶在纤维上为什么表现为 T-对偶。
\par\smallskip\noindent\emph{提示：}环面上的平坦 $U(1)$ 丛（带平坦联络）与对偶环面的点一一对应；KW 用 SYZ 图景作启发，并未把它当作证明。
\end{enumerate}

\subsubsection*{第 5 讲\quad 膜与线算子}

\begin{enumerate}[label=\arabic*.,start=17,leftmargin=2.4em,itemsep=6pt]
\item 判断下列子流形在 $I,J,K$ 下是复子流形还是 Lagrange 子流形，并据此写出膜类型：单点、整个 $\MH$、Hitchin 纤维 $F$、固定 $E$ 而让 $\varphi$ 变化的族、典范余各向同性膜的支撑 $W$。
\par\smallskip\noindent\emph{提示：}点与整个 $\MH$ 在三个复结构下都复（型 $(B,B,B)$）；Hitchin 纤维与固定 $E$ 的族是 $\Omega_I$-Lagrangian 复子流形（型 $(B,A,A)$）；$W$ 在 $J$ 中复、在 $I,K$ 中余各向同性（型 $(A,B,A)$）。

\item 设 $B_p$ 是支撑在 $p\in\MH$ 的 zerobrane，其 Chan--Paton 对象是点 $p$ 上的天穹层。利用特征膜的定义 $XB=B\otimes V$，说明 Wilson 线算子作用在 $B_p$ 上只能把它变成 $B_p\otimes R$，从而 $B_p$ 是电特征膜的原型。
\par\smallskip\noindent\emph{提示：}线的效应只会改变 Chan--Paton 数据；张量一个向量空间不改变膜的支撑，所以特征值不是数而是向量空间。

\item 说明 't Hooft 算子所带的磁荷为什么取值在余权格中，并解释其奇异性如何由三维单极子解实现。
\par\smallskip\noindent\emph{提示：}GNO 量子化条件把磁荷格定为权格的对偶；万有覆盖上的单极子解在远处给出 Dirac 型奇异性。

\item 阅读 KW 第 6.4 节末段，说明为什么电特征膜与磁特征膜都自动带 $C$ 上的平坦向量丛；再解释这个平坦向量丛为什么被解释为 $\Ldual G$-局部系统而不是 $G$-局部系统。
\par\smallskip\noindent\emph{提示：}$\pi_1(C)$ 的表示给出局部系统；Wilson 算子的标签是 $G$ 的表示 $R$，而特征值随 $R$ 变化的规律由对偶群的局部系统编码。
\end{enumerate}

\subsubsection*{第 6 讲\quad 't Hooft 与 Hecke}

\begin{enumerate}[label=\arabic*.,start=21,leftmargin=2.4em,itemsep=6pt]
\item 对 $G=GL(1)$ 写出点 $x$ 处的 Hecke 修改，并证明只有 $L\mapsto L(x)$ 与 $L\mapsto L(-x)$ 两种。
\par\smallskip\noindent\emph{提示：}在 $C\setminus x$ 上同构、在 $x$ 处允许一阶零点或一阶极点；秩 $1$ 时局部数据只由 $\ZZ$ 的一个整数标记。

\item 写出 Hecke 特征层的定义式，并说明它与特征膜定义 $XB=B\otimes V$ 为什么形式相同；解释这个形式上的对应在 KW 词典中意味着什么。
\par\smallskip\noindent\emph{提示：}Hecke 算子把层送到自身张量局部系统的纤维，正如线算子把膜送到自身张量一个向量空间；两者都是「本征值不是数而是空间」的例证。

\item 把四维自对偶方程沿一个平移方向约化，得到三维的 Bogomolny 方程，并说明为什么 't Hooft 算子沿 $C$ 的奇异性对应三维单极子。
\par\smallskip\noindent\emph{提示：}把被约化方向的分量解释为标量 $\varphi$，自对偶条件化为 $F=\ii\star D\varphi$；有限能量解在空间无穷远处的渐近行为就是单极子奇异性。

\item 说明为什么点 $x$ 处的 Hecke 修改由仿射 Grassmannian $G_\CC((z))/G_\CC[[z]]$ 参数化，并说明其 $G_\CC[[z]]$-轨道由支配余权标记。
\par\smallskip\noindent\emph{提示：}在 $x$ 附近取形式坐标，修改的数据是形式圆盘上的一个转移函数；按极点阶数做 Birkhoff 分解即得轨道分解。
\end{enumerate}

\subsubsection*{第 7 讲\quad 翻译器}

\begin{enumerate}[label=\arabic*.,start=25,leftmargin=2.4em,itemsep=6pt]
\item 从典范余各向同性膜（型 $(A,B,A)$）的边界条件出发，解释为什么其开弦代数（$(B_{\mathrm{c.c.}},B_{\mathrm{c.c.}})$ 弦）是作用在 $K_M^{1/2}\otimes L$ 上的微分算子层。
\par\smallskip\noindent\emph{提示：}$J$ 中的 B-型条件给出全纯性，$I,K$ 中的 A-型条件把一半的振荡模消掉；量子化后剩下的态是乘子与微分的组合，即微分算子，$K^{1/2}$ 因子来自测度与反常。

\item 已知 $f$ 是线性的，且对整数 $\Psi=n$ 有 $f(n)=n$。由此证明 $f(\Psi)=\Psi$。
\par\smallskip\noindent\emph{提示：}设 $f(\Psi)=a\Psi+b$，对所有整数 $n$ 代入 $f(n)=n$，比较两个整数点即可。

\item 用 $f(\Psi)=\Psi$ 解释 c.c.\ 膜的弦代数如何随 $\Psi$ 改变；说明为什么 $\Psi=0$ 时它退化为 $K_M^{1/2}$ 上通常 $D$-模的扭化版本。
\par\smallskip\noindent\emph{提示：}$f(0)=0$ 时 $L^{f(\Psi)}=\mathcal O$，只剩 $K_M^{1/2}$ 因子；$\Psi\neq0$ 时多出线丛 $L^\Psi$ 的扭化。
\end{enumerate}

\subsubsection*{第 8 讲\quad 收获与边界}

\begin{enumerate}[label=\arabic*.,start=28,leftmargin=2.4em,itemsep=6pt]
\item 列出 KW 论证链中依赖以下物理假设的步骤：(1) 扭化后的拓扑场论存在且在 $S$-对偶下等价；(2) 膜与 Hecke 特征层的词典；(3) 强耦合与弱耦合之间的解析延拓。说明哪一条最难完全数学化。
\par\smallskip\noindent\emph{提示：}见第 8 讲的总结表；难点集中在把 $S$-对偶写成数学上严格的等价，以及膜范畴与 $D$-模范畴的比较。

\item$^\dagger$ 2024 年 Gaitsgory 等人的工作是纯代数几何的证明。说明它证明了什么、没有证明什么（是否包含 $S$-对偶？是否覆盖分歧情形？），并解释为什么这不妨碍 KW 的物理论证仍然有价值。
\par\smallskip\noindent\emph{提示：}该证明处理特征 $0$、非分歧的 de Rham 情形，证明 Hecke 特征层存在；它不构造物理的 $S$-对偶，也不覆盖分歧与量子化参数的全部范围。

\item 若只考虑非分歧情形，KW 的构造需要哪些额外限制（如局部系统在每点局部平凡、't Hooft 奇异性无冒泡）？要推广到分歧与野分歧，需要补充什么数据？
\par\smallskip\noindent\emph{提示：}分歧情形要在点上加入规范场奇异性的数据（对应抛物结构或 Level 结构），野分歧还需要高阶极点/幂零单值性的信息。
\end{enumerate}
